
\documentclass[oneside,senior,etd]{BYUPhys}

\usepackage{multirow}
\usepackage{longtable}
\usepackage{array}
\usepackage{cmap}
\usepackage{tempora}
\usepackage[T2A]{fontenc}
\usepackage[utf8]{inputenc}
\usepackage[russian]{babel}
\usepackage{DejaVuSansMono}
\usepackage{amsfonts}
\usepackage{amstext}
\usepackage{amssymb}
\usepackage{amsthm}
\usepackage{amsmath}
\usepackage{graphicx}
\usepackage{subfig}
\usepackage{color}
\usepackage[nottoc]{tocbibind}
\usepackage{verbatim}
\usepackage{listings}
\usepackage{tikz}
\usepackage{pgfplots}
\usetikzlibrary{arrows,positioning,shapes,fit,calc}
\usepackage{adjustbox}
\usepackage{makecell}
\usepackage{booktabs}
\usepackage{boldline}
\usepackage{xcolor}
\usepackage{soul}
\usepackage{url}
\usepackage{multirow}
\usepackage{pifont}
\usepackage{indentfirst}
\usepackage{caption}
\usepackage{array}
\usepackage{longtable}
\usepackage{algpseudocode}
\usepackage{placeins}
\usepackage{perpage}
\MakePerPage{footnote}

\definecolor{backgroundColor}{rgb}{0.93, 0.93, 0.93}
\definecolor{commentColor}{rgb}{0.4, 0.4, 0.6}
\definecolor{keyWordsColor}{rgb}{0.4, 0.4, 1}
\definecolor{identifiersColor}{rgb}{0.4, 0.2, 0.6}

\lstdefinelanguage{Kotlin}{
  morekeywords={package,import,class,interface,object,fun,val,var,private,internal,override,return,when,is,if,else,for,in,while,true,false,null,companion,data,sealed,enum,lateinit,by,suspend},
  sensitive=true,
  morecomment=[l]{//},
  morecomment=[s]{/*}{*/},
  morestring=[b]",
}

\lstset{
  language=Kotlin,
  basicstyle=\footnotesize\ttfamily,
  numbers=left,
  numberstyle=\tiny,
  numbersep=5pt,
  tabsize=2,
  extendedchars=true,
  breaklines=true,
  frame=b,
  stringstyle=\color{blue}\ttfamily,
  showspaces=false,
  showtabs=false,
  xleftmargin=17pt,
  framexleftmargin=17pt,
  framexrightmargin=5pt,
  framexbottommargin=4pt,
  commentstyle=\color{commentColor},
  keywordstyle=\color{keyWordsColor},
  identifierstyle=\color{identifiersColor},
  backgroundcolor=\color{backgroundColor},
  escapeinside={(*}{*)}
}

\definecolor{importantTextColor}{rgb}{0.8, 0.2, 0.2}
\newcommand{\highlight}[1]{\textcolor{importantTextColor}{\fontsize{12}{11}\ttfamily\linespread{1.0}#1}}
\newcommand{\code}[1]{\lstinline|#1|}
\usepackage[table]{colortbl}

\Faculty{Факультет вычислительной математики и кибернетики}
\Chair{Кафедра алгоритмических языков}
\Year{2026}
\Month{Май}
\City{Москва}
\AuthorText{Автор:}
\Author{Плотников Алексей Сергеевич}
\AuthorEng{Student Name}
\AcadGroup{000}

\TitleTop{Сравнение моделей машинного обучения}
\TitleMiddle{для поведенческой аутентификации}
\TitleBottom{в библиотеке андроид-приложения}

\docname{Курсовая работа}
\Advisor{Полякова Ирина Николаевна}
\AdvisorDegree{к.ф.-м.н.}
%\Consultant{}
%\ConsultantDegree{}

\Abstract{
Работа посвящена задаче непрерывной поведенческой аутентификации пользователя мобильного устройства и сравнению моделей машинного обучения в составе библиотеки андроид-приложения. Отправной точкой является проблема одноразовой аутентификации: пароль, пин-код и отпечаток пальца проверяют пользователя в момент входа, но не гарантируют, что после разблокировки устройством продолжает пользоваться тот же человек.

Поведенческая аутентификация состоит из запоминания обычного поведения пользователя, анализа временных окон поведения и последующего дообучения при отсутствии признаков смены пользователя. В работе проводится сравнение комбинаций моделей: детекторов аномальности временных окон и детекторов смены пользователя. Для принятия решения об аномальности временного окна поведения рассматриваются статистический подход, MiniBatch K-Means, модель гауссовых смесей, автоэнкодер и изолирующий лес; для принятия решения о смене пользователя рассматриваются эвристическое голосование, скрытая марковская модель и байесовское обнаружение точек изменения.

Программной основой работы является библиотека для андроид, реализованная на Kotlin, собирающая данные касаний, инерциальных сенсоров и контекстные признаки, преобразующая поток событий во временные окна и передающая нормализованные признаки моделям. Отдельно описана запись сырых событий в построчном JSONL-формате и последующее воспроизведение тех же событий для сравнения моделей на одном наборе данных. В работе предложена методика сбора данных через встраивание библиотеки в приложения. Для сравнения моделей используются не только показатели FAR, FRR, EER и Recall, но и задержка обнаружения смены пользователя, нагрузка на процессор, оперативную память и батарею.
}

\TitlePageText{Титульная страница}
\University{Московский государственный университет имени М.В.Ломоносова}
\GrText{гр.}
\AdvisorText{Научный руководитель}
\ConsultantText{Научный консультант}
\AbstractText{Аннотация}
\AcknowledgmentsText{Благодарности}
\ListingText{Листинг}
\AlgorithmText{Алгоритм}

\hypersetup{
  pdftitle={\PDFTitle},
  pdfauthor={\PDFAuthor},
  bookmarks=false
}

\begin{document}
\emergencystretch=3em
\sloppy

\fixmargins
\makepreliminarypages
\oneandhalfspace

% PDF bookmarks disabled for stable pdfLaTeX build
\tableofcontents
\clearpage

\section{Введение}

Современный смартфон является не только средством связи, но и постоянной точкой доступа к личным и рабочим данным пользователя. Через мобильное устройство открываются банковские приложения, корпоративная почта, мессенджеры, медицинские сервисы, облачные хранилища, документы и учетные записи в социальных сетях. Поэтому для мобильной платформы важна не только защита момента входа, но и контроль всей дальнейшей пользовательской сессии. Если телефон разблокирован и передан другому человеку, оставлен на столе, украден из рук или используется в присутствии посторонних, классические механизмы входа уже не решают задачу подтверждения текущего пользователя.

Существует множество методов аутентификации. В мобильных приложениях чаще всего применяются пароль, пин-код, графический ключ, отпечаток пальца или распознавание лица. Эти способы удобны и хорошо изучены, однако их общая слабость состоит в том, что проверка выполняется в конкретные моменты времени: при входе или при отдельной критической операции. Между такими проверками система обычно предполагает, что устройство все еще находится у легитимного пользователя.

Поведенческая аутентификация решает эту проблему иначе. Система не спрашивает пользователя о секрете и не требует отдельного действия. Она наблюдает за уже возникающим поведением: как пользователь держит устройство, как нажимает на экран, с какой скоростью выполняет свайпы, как изменяется ориентация телефона во время работы, в какие моменты возникают паузы и какова ритмика взаимодействия. Если текущее поведение значительно отличается от профиля владельца, система может зафиксировать возможную смену пользователя и вызвать дополнительную проверку. Такой механизм не заменяет пароль или биометрический сканер, а дополняет их, закрывая промежуток между первоначальной аутентификацией и завершением сессии.

\subsection{Проблема одноразовой аутентификации}

Пароль, пин-код и отпечаток пальца проверяют пользователя в конкретный момент времени. После успешной проверки приложение или операционная система обычно открывает доступ на некоторый период. Этот период может быть коротким, например при подтверждении банковской операции, или длинным, например при работе в почтовом клиенте или корпоративном мессенджере. В обоих случаях возникает окно риска: приложение продолжает доверять устройству, хотя фактический человек перед экраном мог измениться.

Такая ситуация особенно характерна для смартфонов. В отличие от стационарного компьютера, телефон часто используется на ходу, в транспорте, в аудитории, дома или на рабочем месте. Пользователь может временно передать устройство знакомому для просмотра фотографии, ребенку для игры, коллеге для чтения сообщения. Устройство может быть оставлено разблокированным или украдено уже после разблокировки. Если приложение полагается только на начальный вход, оно не имеет механизма обнаружения такой смены пользователя.

Непрерывная аутентификация вводит дополнительный уровень контроля. В ней проверка не является разовым событием, а повторяется неявно, по мере поступления новых данных. При этом система не должна превращать безопасность в постоянное раздражение пользователя. Поведенческие признаки ценны именно тем, что их можно собирать пассивно. Пользователь не вводит специальную фразу, не проходит тест и не выполняет дополнительный жест: он просто использует приложение, а библиотека анализирует естественные действия.

\subsection{Цель и задачи работы}

Цель работы состоит в разработке и описании подхода к сравнению моделей машинного обучения для поведенческой аутентификации пользователя в библиотеке андроид-приложения. Под сравнением понимается не только оценка точности моделей, но и оценка пригодности для мобильной среды: задержки принятия решения, загрузки процессора, расхода памяти и энергии.

Для достижения цели требуется решить следующие задачи:
\begin{enumerate}
    \item Проанализировать ограничения традиционной аутентификации и место поведенческой аутентификации в защите пользовательской сессии;
    \item Рассмотреть существующие проблемы исследований поведенческой аутентификации;
    \item Определить набор поведенческих признаков, доступных на андроид-устройстве;
    \item Спроектировать конвейер, состоящий из сбора сырых событий, построения временных окон, оценки аномальности окон и анализа последовательности аномалий;
    \item Реализовать архитектуру библиотеки, в которой предусмотрено расширение набора признаков и подключение новых моделей;
    \item Описать и сравнить детекторы аномальности временных окон;
    \item Описать и сравнить модели определения смены пользователя;
    \item Сравнить пары моделей обнаружения аномалий и определения смены пользователя по метрикам FAR, FRR, EER, Recall, задержке обнаружения, дополнительной нагрузке на процессор, дополнительному использованию оперативной памяти и энергопотреблению.
\end{enumerate}


\section{Обзор существующих исследований и проблем существующих подходов}

\subsection{Исследования по клавиатурному, мышиному и мобильному поведению пользователя}
Рассмотрим работы в смежной области. Например в работе Р.Д.~Заклякова, М.И.~Петровского и И.С.~Попова рассматривалась разработка гибридных методов распознавания пользователя по особенностям работы со стандартными устройствами ввода компьютера~\cite{zaklyakov2017hybrid}. В качестве источников поведенческой информации в такой постановке выступают клавиатура и мышь. Пользователь не предъявляет отдельный биометрический образец, а проявляет индивидуальные особенности в обычных действиях: нажимает клавиши с характерными задержками, перемещает указатель с определенной скоростью, выполняет последовательности операций в своем темпе.

Методологическая ценность работы состоит в использовании нескольких источников информации и в рассмотрении гибридного распознавания. Для поведенческой аутентификации это особенно важно, поскольку один источник ввода часто оказывается неполным. Например, клавиатурный почерк информативен при наборе текста, но почти не помогает при работе с графическим интерфейсом; мышиная динамика отражает особенности управления курсором, но не описывает ритм набора. Объединение таких признаков повышает устойчивость распознавания и показывает общий принцип: поведение пользователя следует рассматривать не как отдельное действие, а как совокупность повторяющихся микродвижений и временных характеристик.

Ограничение этой линии исследований связано с предметной областью. Стандартные устройства ввода персонального компьютера работают в относительно стабильной среде. Клавиатура и мышь обычно находятся на столе, положение руки меняется ограниченно, а поток событий имеет понятную структуру. Смартфон используется иначе: устройство перемещается в руке, может находиться под разными углами, взаимодействие состоит из касаний, свайпов, удержаний, поворотов и часто зависит от типа приложения. Поэтому идеи, разработанные для клавиатуры и мыши, нельзя механически перенести в андроид-приложение. Их следует рассматривать как доказательство информативности поведенческих признаков и как методологическую основу для перехода к более сложной мобильной среде.

Вторая важная работа --- исследование Е.И.~Краснякова и И.С.~Попова, посвященное аутентификации пользователя мобильного устройства по характеристикам работы с виртуальной клавиатурой~\cite{krasnyakov2017virtualkeyboard}. Эта работа ближе к теме работы, поскольку переносит анализ поведения на мобильное устройство и рассматривает экранный ввод как источник индивидуальных признаков. Виртуальная клавиатура отличается от физической: пользователь нажимает не механические клавиши, а области сенсорного экрана; длительность касания, точность попадания, паузы между символами и характер исправления ошибок могут зависеть от манеры держать телефон и от привычек пользователя. Поэтому виртуальная клавиатура является естественным мостом между классическим клавиатурным почерком и современной мобильной поведенческой биометрией.

Однако аутентификация по виртуальной клавиатуре покрывает только часть поведения пользователя. В современных приложениях значительная часть взаимодействия не сводится к набору текста. Пользователь читает, прокручивает, выбирает элементы интерфейса, удерживает телефон во время ожидания, меняет положение устройства, играет, выполняет быстрые жесты или почти не касается экрана. Если система опирается только на ввод текста, она получает данные нерегулярно и может долго не иметь достаточного основания для решения. Это особенно заметно в приложениях для чтения, просмотра документов или пассивного потребления контента. Поэтому для библиотеки поведенческой аутентификации необходимо расширить область наблюдения: учитывать не только клавиатурные события, но и общие касания, жесты и сенсорный фон.

В этих работах также важна общая исследовательская позиция: поведенческие признаки не являются абсолютно неизменными, как некоторые физиологические биометрические характеристики. Пользователь может устать, поменять руку, изменить темп работы, использовать другое устройство или оказаться в другом контексте. Это означает, что модель должна быть не только распознающей, но и устойчивой к естественной вариативности поведения. Для работы отсюда следует два вывода. Во-первых, систему нельзя оценивать только на одном коротком сценарии, потому что такой сценарий может переоценивать качество. Во-вторых, для сравнения моделей важно фиксировать протокол: если разные модели обучаются и проверяются на разных фрагментах поведения, то различия в результатах будут объясняться не только алгоритмами.

\subsection{Современные обзорные работы по мобильной поведенческой биометрии}

Обзорные работы показывают, что мобильная поведенческая аутентификация стала самостоятельной и достаточно широкой областью исследований. В обзоре I.~Stylios, S.~Kokolakis, O.~Thanou и S.~Chatzis систематизируются технологии поведенческой биометрии и непрерывной аутентификации на мобильных устройствах~\cite{stylios2021survey}. Важный вывод этого обзора состоит в том, что мобильная аутентификация не ограничивается одной модальностью. В литературе рассматриваются клавиатурная динамика, жесты на сенсорном экране, движение устройства, походка, голос, особенности использования приложений и их комбинации. Такой обзор полезен для настоящей работы потому, что показывает: задача выбора модели не может быть отделена от выбора источников данных и способа их предварительной обработки.

Одна из проблем, отмечаемых в обзорных работах, состоит в неоднородности экспериментальных условий. Исследования используют разные наборы пользователей, разные устройства, разные сценарии, разные частоты сенсоров и разные способы выделения признаков. В результате итоговые значения точности, FAR, FRR или EER часто нельзя напрямую сравнивать между публикациями. Например, модель, проверенная на контролируемом наборе жестов, может показывать лучшие числа, чем модель, проверенная на свободном взаимодействии с приложением, но это не означает ее превосходства в реальной среде. Необходимо сравнивать модели на одном и том же потоке данных, а не переносить оценки из разных статей.

В обзоре M.~Abuhamad, A.~Abusnaina, D.~Nyang и D.~Mohaisen рассматривается непрерывная аутентификация пользователей смартфонов на основе данных датчиков и анализируется большое количество работ, посвящённых подходам по движениям устройства, походке, динамике набора текста, касаниям экрана, голосовым признакам и сочетанию нескольких видов поведенческих данных.~\cite{abuhamad2021contemporary}. Эта классификация показывает, что мобильный телефон предоставляет множество источников неявной информации о пользователе. При этом разные группы методов имеют разные ограничения. Методы, основанные на походке, хорошо работают в движении, но менее полезны, когда пользователь сидит. Методы, основанные на касаниях, требуют активного взаимодействия с экраном. Методы, основанные на сенсорах, могут получать данные чаще, но чувствительны к положению устройства и внешнему контексту. Следовательно, практическая андроид-библиотека должна быть готова к неполному и неравномерному потоку признаков.

Обзор A.~Ray-Dowling и соавторов, посвященный стационарной мобильной поведенческой биометрии, дополняет эту картину тем, что обращает внимание на сценарии, в которых пользователь не обязательно идет или активно перемещается~\cite{raydowling2023stationary}. Для мобильной аутентификации это важно: смартфон часто используется в положении сидя, лежа, стоя в транспорте или во время кратких взаимодействий. В таких условиях походка не является основным источником данных, а более значимыми становятся микродвижения устройства, касания, удержание телефона и ритм работы с интерфейсом. Это особенно близко к рассматриваемой библиотеке, поскольку она должна работать в приложениях разных типов, включая пассивное чтение, набор текста и активное взаимодействие.

Обзорные работы также поднимают вопрос о балансе между безопасностью и удобством. Слишком чувствительная модель будет часто ошибочно отвергать владельца, создавая ложные срабатывания и снижая доверие к системе. Слишком мягкая модель будет пропускать злоумышленника. Поэтому качество поведенческой аутентификации нельзя описать одной метрикой accuracy. Необходимо рассматривать ошибки разных типов, задержку до принятия решения и практическую стоимость ложного предупреждения. Этот вывод далее используется в работе при формировании набора метрик сравнения моделей, но в рамках текущего раздела важно подчеркнуть именно происхождение проблемы: она связана с самой природой непрерывной аутентификации, где решение принимается многократно и в условиях шума.

Еще одна проблема состоит в слабой воспроизводимости результатов. Даже если авторы подробно описывают признаки и модель, повторить эксперимент бывает трудно из-за закрытого набора данных, отсутствия исходного кода, различий в сенсорах или недостаточного описания разбиения на обучение и тест. Модель должна оцениваться в условиях, при которых алгоритм обработки изменяется, а исходный поток событий остаётся неизменным.

\subsection{Открытые базы данных и соревнования как попытка стандартизировать сравнение}

Часть проблем сравнения решается за счет открытых баз данных и соревнований. Одним из наиболее показательных примеров является MobileB2C --- соревнование по мобильной поведенческой биометрии, описанное G.~Stragapede и соавторами~\cite{mobileb2c2022}. Его цель состояла в сравнении систем аутентификации, использующих поведенческие признаки, которые прозрачно собираются мобильными устройствами при обычном взаимодействии человека с интерфейсом. В соревновании выделялись несколько типичных активностей: набор текста, чтение текста, свайпы в галерее и касания. Такой выбор сценариев важен, поскольку демонстрирует: поведение пользователя нельзя оценивать только в одном режиме. В разных приложениях возникают разные типы событий, и система должна быть проверена на нескольких видах активности.

MobileB2C особенно важен для работы как пример корректной постановки сравнения. Несколько участников работают с общей базой и общим протоколом, поэтому различия в результатах в большей степени отражают различия в алгоритмах. Это лучше, чем сравнение статей, где каждая статья использует собственный набор данных. Тем не менее формат соревнования имеет и ограничения. Он стандартизирует исследовательскую оценку, но не обязательно решает задачу встраивания в реальное приложение. В соревновании данные уже собраны, а библиотеке необходимо дополнительно решить вопросы получения событий, работы в фоне, хранения, воспроизведения, конфигурации, нагрузки на устройство и интеграции с приложением.

С близкой целью была создана база BehavePassDB, представленная G.~Stragapede, R.~Vera-Rodriguez, R.~Tolosana и A.~Morales~\cite{behavepassdb2023}. База ориентирована на мобильную поведенческую биометрию и включает типичные активности взаимодействия с устройством. Важное достоинство этой работы --- акцент на разделении пользовательского и устройственного факторов. В мобильной биометрии существует риск, что модель обучается отличать не столько пользователя, сколько конкретное устройство, частоту сенсоров, размер экрана или особенности регистрации событий. Если такой эффект не контролировать, качество на тесте может оказаться завышенным, а при переносе на другое устройство модель будет работать хуже.

Для настоящей работы BehavePassDB важна не как готовая замена собственного сбора данных, а как пример требований к протоколу. Из нее следует, что при проектировании эксперимента нужно учитывать несколько сессий, разные задачи, возможность работы разных пользователей с одним устройством и различение сценариев случайного и более подготовленного нарушителя. Встраиваемая библиотека должна обеспечивать сбор данных в формате, пригодном для последующего воспроизводимого сравнения. Если же данные сохраняются только в виде уже агрегированных признаков одной версии алгоритма, то последующая проверка новых гипотез становится затруднительной.

Открытые базы данных также показывают, что не существует универсального числа качества, которое полностью описывает систему. Один алгоритм может лучше работать на наборе текста, другой --- на свайпах, третий --- на сенсорных данных. Кроме того, метод может показывать хорошее среднее качество, но плохо работать на части пользователей. Для библиотеки, предназначенной для реальных приложений, это означает необходимость не только выбрать одну лучшую модель, но и иметь экспериментальный контур, где можно сравнивать несколько моделей и несколько способов принятия решения на одинаковых данных.

\subsection{Примеры конкретных систем и моделей непрерывной аутентификации}

Наряду с обзорными работами и базами данных существуют исследования, предлагающие конкретные модели непрерывной аутентификации. Одним из таких примеров является AuthConFormer, предложенный M.~Hu, K.~Zhang, R.~You и B.~Tu~\cite{hu2023authconformer}. Авторы рассматривают сенсорную непрерывную аутентификацию пользователей смартфона и предлагают архитектуру на основе сверточного трансформера. Существенная особенность этой работы состоит в попытке совместить способность глубоких моделей извлекать сложные признаки из сырых сенсорных данных с требованием пригодности для мобильной платформы. В статье отдельно обсуждается возможность работы на процессорах мобильных устройств, что отличает ее от исследований, где модель оценивается только с точки зрения точности.

AuthConFormer демонстрирует достоинства современных глубоких подходов: они способны извлекать более сложные зависимости из временных рядов датчиков и потенциально лучше работать с нелинейными паттернами поведения. Однако одновременно эта работа показывает и проблему. Сложная модель требует большего объема данных, более аккуратного обучения, настройки архитектуры и контроля вычислительной стоимости. Для библиотеки, которая должна поддерживать несколько вариантов моделей, это означает, что глубокую модель нельзя рассматривать изолированно. Ее необходимо сравнивать с более простыми методами на тех же данных и с учетом ресурсов устройства. Иначе существует риск выбрать алгоритм, который выглядит лучшим на исследовательском стенде, но оказывается избыточным или нестабильным в приложении.

Другой пример --- работы, развивающие системы аутентификации по касаниям и нескольким контекстам взаимодействия. В статье P.M.A.B.~Estrela и соавторов предлагается развитие Biotouch Framework для непрерывной аутентификации на основе обучения с учителем~\cite{estrela2021framework}. Важная для настоящей работы идея состоит в том, что поведение пользователя следует анализировать не в одном общем режиме, а с учетом разных областей или scopes взаимодействия. Это согласуется с практическим наблюдением: одно и то же действие в разных частях приложения может иметь разный смысл, а признаки касаний зависят от интерфейса и пользовательской задачи. Ограничение такого подхода состоит в необходимости заранее определить сценарии и собрать достаточное число примеров для каждого из них.

В литературе также широко обсуждаются одноклассовые и аномалийные постановки. Они особенно естественны для мобильной защиты, поскольку данные злоумышленников заранее неизвестны. Однако эта постановка сложнее, чем обычная бинарная классификация. Модель обучается на поведении владельца и должна отделять нормальные отклонения владельца от поведения другого человека. Если профиль слишком узкий, он будет часто выдавать ложные тревоги; если слишком широкий, он может принять нарушителя. Поэтому в исследованиях используются разные подходы: статистические модели, метод опорных векторов для одного класса, автоэнкодеры, деревья изоляции, кластеризация, вероятностные модели. Обзорные работы показывают многообразие таких методов, но одновременно подчеркивают отсутствие единого победителя для всех условий~\cite{stylios2021survey,abuhamad2021contemporary}.

Сравнение конкретных моделей в опубликованных работах часто ограничено рамками выбранной статьи. Авторы предлагают новый алгоритм и сравнивают его с несколькими базовыми решениями. Это естественно для научной публикации, но для инженерной библиотеки такого сравнения недостаточно. В приложении важна не только максимальная точность, но и стабильность при изменении сценария, возможность локального обучения, скорость получения решения, объем хранимого профиля, простота добавления новых признаков и интерпретируемость предупреждения. Поэтому настоящая работа опирается на опубликованные модели как на набор кандидатов и идей, но формулирует отдельную задачу: построить среду, в которой такие модели проверяются по единому протоколу.

\subsection{Сводное сравнение рассмотренных работ}

Для удобства сопоставления рассмотренные исследования сведены в таблицу~\ref{tab:related-work}. Таблица не повторяет техническое описание разрабатываемой библиотеки, а показывает, какие идеи из существующих работ важны для постановки работы и какие ограничения остаются после их анализа.

\begin{longtable}{|p{0.24\textwidth}|p{0.30\textwidth}|p{0.34\textwidth}|}
\caption{Сравнение существующих работ по поведенческой аутентификации}\label{tab:related-work}\\
\hline
\textbf{Работа} & \textbf{Что рассматривается} & \textbf{Ограничение для задачи} \\
\hline
\endfirsthead
\hline
\textbf{Работа} & \textbf{Что рассматривается} & \textbf{Ограничение для задачи} \\
\hline
\endhead
Закляков, Петровский, Попов~\cite{zaklyakov2017hybrid} & Распознавание пользователя по клавиатуре и мыши; гибридное использование стандартных устройств ввода & Работа относится к персональному компьютеру; мобильные сенсоры, жесты и условия андроид-приложения не являются основным объектом исследования \\
\hline
Красняков, Попов~\cite{krasnyakov2017virtualkeyboard} & Аутентификация пользователя мобильного устройства по виртуальной клавиатуре & Рассматривается важная, но узкая часть мобильного поведения; не покрываются пассивное использование, общие жесты и сенсорный фон \\
\hline
Stylios и соавт.~\cite{stylios2021survey} & Обзор поведенческой биометрии и непрерывной аутентификации на мобильных устройствах & Систематизирует область, но не дает единого экспериментального контура для сравнения моделей в библиотеке \\
\hline
Abuhamad и соавт.~\cite{abuhamad2021contemporary} & Обзор сенсорной непрерывной аутентификации; классификация методов по типам поведенческих данных & Показывает разнообразие модальностей, но результаты отдельных работ остаются трудно сопоставимыми из-за разных данных и протоколов \\
\hline
MobileB2C~\cite{mobileb2c2022} & Соревнование и общий протокол для сравнения мобильных систем поведенческой биометрии & Улучшает честность сравнения, но не решает полностью вопросы интеграции, записи событий и работы библиотеки внутри приложения \\
\hline
BehavePassDB~\cite{behavepassdb2023} & Публичная база данных с несколькими мобильными HCI-задачами и учетом device bias & Полезна как эталон протокола, но не заменяет механизм сбора и воспроизведения данных в конкретном приложении \\
\hline
AuthConFormer~\cite{hu2023authconformer} & Глубокая модель на основе сверточного трансформера для сенсорной непрерывной аутентификации & Перспективна по качеству, но требует сравнения с более простыми моделями по ресурсам, данным и устойчивости \\
\hline
\end{longtable}

Из таблицы видно, что существующие работы решают разные части общей задачи. Одни доказывают информативность поведенческих признаков, другие систематизируют модальности, третьи предлагают базы и протоколы, четвертые строят конкретные модели. Однако для темы работы важна их совокупность, а не отдельная публикация. Поэтому дальнейшая постановка работы строится как объединение этих требований.

\subsection{Выводы по обзору и нерешенные проблемы}

Рассмотренные исследования позволяют выделить несколько нерешенных проблем. Первая проблема --- фрагментарность признаков. Значительная часть работ анализирует один источник данных: клавиатуру, касания, свайпы, движение или походку. Такой подход удобен для контролируемого эксперимента, но в реальном мобильном приложении пользовательское поведение неоднородно. В одних сценариях много текстового ввода, в других --- жестов, в третьих --- почти только удержание устройства. Поэтому система, ориентированная на практическое использование, должна поддерживать несколько групп признаков и сохранять работоспособность при временном отсутствии части событий.

Вторая проблема --- несопоставимость опубликованных результатов. Даже современные обзоры не могут превратить результаты разных работ в полностью корректный рейтинг моделей, потому что статьи используют разные базы, протоколы и метрики. Из этого следует требование к работе: сравнение моделей должно выполняться на одном наборе данных и при одинаковом способе построения признаков. Только в этом случае можно обсуждать различия между статистическим подходом, кластеризацией, вероятностной моделью, автоэнкодером или деревьями изоляции как различия между алгоритмами, а не как следствие разных условий эксперимента.

Третья проблема --- недостаточное внимание к системным ограничениям мобильного устройства. В публикациях качество модели часто оказывается главным результатом, тогда как нагрузка на процессор, память и батарею рассматривается реже. Между тем непрерывная аутентификация должна работать в фоне и не мешать пользователю. Поэтому тяжелая модель с небольшим выигрышем в точности может быть хуже простой модели, если она быстрее разряжает устройство или увеличивает задержки интерфейса. Для библиотеки это не второстепенная деталь, а часть критерия пригодности.

Четвертая проблема --- отделенность исследовательских решений от интеграции в реальные приложения. Научная статья может описать модель и эксперимент, но не обязательно предоставляет механизм, позволяющий разработчику подключить поведенческую аутентификацию к приложению, собрать события, сохранить их, воспроизвести и заменить модель. Для настоящей работы это принципиально: объектом является не только алгоритм, но и библиотека, в которой алгоритмы должны быть взаимозаменяемыми.

Таким образом, обзор существующих исследований приводит к следующей постановке. Необходимо рассматривать поведенческую аутентификацию как практическую мобильную систему, в которой используются идеи клавиатурного и мышиного поведения, мобильной виртуальной клавиатуры, сенсорной биометрии, открытых протоколов сравнения и современных моделей машинного обучения. При этом основной исследовательский акцент должен быть сделан на воспроизводимом сравнении подходов: одинаковый поток данных, одинаковое построение признаков, несколько моделей, несколько критериев качества и учет ограничений устройства.


\section{Предлагаемый подход}

\subsection{Требования к библиотеке}

Из описанных проблем следуют требования к системе. Библиотека должна собирать широкий спектр признаков, работать без подключения к интернету, поддерживать разные модели, обеспечивать сопоставимое сравнение моделей на одном наборе данных и не создавать существенной нагрузки на устройство. Кроме того, библиотека должна быть достаточно простой для внедрения в приложение: разработчику не следует вручную обрабатывать каждый сенсор и каждое касание.

\begin{table}[ht]
\centering
\caption{Проблемы существующих подходов и требования к решению}
\label{tab:problem-requirement}
\begin{tabular}{|p{0.42\textwidth}|p{0.46\textwidth}|}
\hline
\textbf{Проблема} & \textbf{Требование к библиотеке} \\
\hline
Проверка пользователя только в момент входа & Непрерывный анализ поведения во время сессии \\
\hline
Малое число признаков & Поддержка касаний, свайпов, сенсоров и контекстных признаков \\
\hline
Сравнение моделей на разных данных & Запись сырых событий и воспроизведение одного и того же потока \\
\hline
Оценка только точности & Учет FAR, FRR, EER, Recall, задержки, дополнительной нагрузки на процессор, дополнительного использования оперативной памяти и расхода батареи \\
\hline
Зависимость от сервера & Локальная обработка на устройстве без обязательного интернета \\
\hline
Сложность расширения & Единые интерфейсы для добавления новых признаков и моделей \\
\hline
Сложность интеграции & Оформление решения как библиотеки с компактным API \\
\hline
\end{tabular}
\end{table}

\subsection{Преимущества предлагаемого подхода}

Предлагаемый подход строится вокруг идеи воспроизводимого экспериментального конвейера. Сначала библиотека записывает поток сырых событий в том виде, в котором они возникают при использовании приложения. Затем тот же поток воспроизводится многократно, меняя длину окна, шаг окна, набор признаков, модель аномальности и модель смены пользователя. Это позволяет сравнивать модели честно: входные данные одинаковы, меняется только алгоритмическая часть.

Второе преимущество состоит в разделении уровней. Детектор аномальности отвечает только на вопрос, похоже ли текущее окно поведения на нормальный профиль владельца. Детектор смены пользователя анализирует последовательность таких ответов. Благодаря этому можно составлять комбинации моделей. Например, статистический детектор можно использовать с эвристическим голосованием, скрытой марковской моделью или байесовским обнаружением точек изменения. То же относится к модели гауссовых смесей, автоэнкодеру и изолирующему лесу.

Третье преимущество --- расширяемость. Предусмотрено добавление новых признаков в слой построения окон, а новые модели --- через соответствующие интерфейсы. При этом остальная часть библиотеки не должна знать внутреннее устройство алгоритма. Такой подход важен для исследовательской работы, где набор моделей и признаки могут меняться.

Четвертое преимущество --- локальность. Все основные этапы выполняются на устройстве. Это уменьшает зависимость от сети и снижает риск утечки данных. Для снижения риска при практическом внедрении целесообразно предусмотреть шифрование профиля и более строгие политики хранения сырых событий.

Система состоит из четырех основных этапов: сбор поведенческой информации, построение временных окон, анализ аномальности окон и анализ последовательности аномалий для определения смены пользователя. Эта схема отражает естественный поток данных от андроид-событий к итоговому решению.

\begin{figure}[ht]
\centering
\begin{tikzpicture}[node distance=1.05cm,>=latex]
\tikzstyle{block}=[rectangle, rounded corners, draw=black, align=center, text width=5cm, minimum height=1.05cm]
\node[block] (collect) {Сбор событий\\касания, сенсоры, контекст};
\node[block, below=of collect] (raw) {Сырые события\\JSONL-запись};
\node[block, below=of raw] (window) {Временные окна\\агрегация признаков};
\node[block, below=of window] (anom) {Детектор\\аномальности окна};
\node[block, below=of anom] (change) {Детектор\\смены пользователя};
\draw[->] (collect) -- (raw);
\draw[->] (raw) -- (window);
\draw[->] (window) -- (anom);
\draw[->] (anom) -- (change);
\node[block, right=2.0cm of raw] (replay) {Воспроизведение\\того же потока};
\draw[->] (raw) -- (replay);
\draw[->] (replay) |- (window);
\end{tikzpicture}
\caption{Конвейер сбора, хранения, обработки и воспроизведения поведенческих данных}
\label{fig:pipeline}
\end{figure}

\clearpage

\section{Сбор поведенческой информации}

\subsection{Источники данных на устройстве}

Платформа андроид предоставляет несколько источников данных, полезных для поведенческой аутентификации. Первый источник --- события касаний. К ним относятся координаты, действие, количество указателей, давление, размер области касания, история движения и технические параметры события~\cite{androidmotionevent}. На основе этих данных вычисляются продолжительность нажатия, скорость свайпа, угол траектории, силу нажатия, площадь контакта и время между нажатиями.

Второй источник --- сенсоры движения и положения. В разработанной библиотеке используются акселерометр, гироскоп, магнитометр, датчик гравитации и линейного ускорения. Средства андроид для работы с датчиками позволяют получать события от доступных сенсоров и регистрировать обработчики изменений~\cite{androidsensors,androidsensorevent}. Эти данные описывают скорость перемещения устройства в руке, скорость вращения вокруг осей и наклон устройства в пространстве.

Третий источник --- контекстные признаки. В разработанной библиотеке используются текущее время и тип сети: Wi-Fi, мобильный интернет или отсутствие подключения. Эти признаки не являются биометрическими в строгом смысле, но помогают описывать обычные условия использования. Например, пользователь может чаще работать с приложением вечером или преимущественно в домашней Wi-Fi-сети. Такие признаки должны использоваться осторожно: они не должны сами по себе определять личность, но могут дополнять моторные признаки.

\subsection{Рассматриваемые поведенческие признаки}

Список признаков в работе строится вокруг тем, которые важны для мобильной поведенческой аутентификации. Признаки разделяются на группы: признаки касания, признаки свайпа, сенсорные признаки и контекстные признаки. В таблице~2 приведена развернутая классификация.

\clearpage

\begin{longtable}{|p{0.12\textwidth}|p{0.3\textwidth}|p{0.5\textwidth}|}
	\caption{Рассматриваемые поведенческие признаки}\label{tab:features-detailed}\\
	\hline
	\textbf{Группа} & \textbf{Признак} & \textbf{Смысл для аутентификации} \\
	\hline
	\endfirsthead
	
	\hline
	\textbf{Группа} & \textbf{Признак} & \textbf{Смысл для аутентификации} \\
	\hline
	\endhead
	
	\multirow{4}{0.2\textwidth}{Касания}
	& Продолжительность нажатия
	& Отражает ритм и моторную привычку взаимодействия с экраном \\
	\cline{2-3}
	& Сила нажатия
	& Может отличаться у пользователей из-за манеры держать устройство и нажимать на экран \\
	\cline{2-3}
	& Площадь касания
	& Косвенно связана с положением пальца, углом касания и размером области контакта \\
	\cline{2-3}
	& Время между нажатиями
	& Описывает темп работы в приложении и паузы между действиями \\
	\hline
	
	\multirow{2}{0.2\textwidth}{Свайпы}
	& Скорость свайпа
	& Характеризует интенсивность жестов и стиль перемещения по экрану \\
	\cline{2-3}
	& Угол траектории
	& Различает вертикальные, горизонтальные и диагональные жесты, а также индивидуальные отклонения \\
	\hline
	
	\multirow{3}{0.2\textwidth}{Сенсоры}
	& Скорость перемещения устройства
	& Вычисляется по акселерометру и линейному ускорению, отражает микродвижения телефона в руке \\
	\cline{2-3}
	& Скорость вращения вокруг осей
	& Измеряется гироскопом и описывает повороты устройства при взаимодействии \\
	\cline{2-3}
	& Наклон устройства
	& Определяется через гравитацию и магнитометр, характеризует привычную позу удержания \\
	\hline
	
	\multirow{2}{0.2\textwidth}{Контекст}
	& Час использования
	& Описывает типичный временной контекст, но должен иметь вспомогательное значение \\
	\cline{2-3}
	& Тип сети
	& Позволяет учитывать обычные условия использования: мобильная сеть, Wi-Fi или отсутствие подключения \\
	\hline
	
\end{longtable}

\subsection{Сбор событий касаний}

Каждое событие преобразуется во внутренний объект, включающий время записи, тип действия, индекс активного указателя, количество указателей, время события, время начала жеста и список указателей. Для каждого указателя сохраняются координаты, исходные координаты, давление, размер, параметры области касания, ориентация и история движения. Такой уровень детализации важен для последующего пересчета признаков: если в будущем потребуется новый признак, его извлечение возможно из сохранённого события.

С точки зрения аутентификации отдельное касание не является устойчивым биометрическим шаблоном. Пользователь может случайно нажать сильнее или задержать палец дольше. Поэтому касания агрегируются. Например, для завершённого жеста вычисляются длительность, пройденное расстояние, скорость, среднее давление и средний размер касания. Затем несколько жестов попадают во временное окно, и по ним строятся усредненные, минимальные и максимальные характеристики.

\subsection{Сбор сенсорных и контекстных данных}

Сенсорные события отличаются от касаний тем, что поступают с собственной частотой, зависящей от устройства, выбранного режима и нагрузки системы. Поэтому нельзя напрямую сравнивать последовательности отдельных измерений, полученных на разных устройствах или в разные моменты. Необходимо агрегировать их во временные окна. В проекте для каждого окна используются значения акселерометра, гироскопа, магнитометра, гравитации и линейного ускорения по осям $x$, $y$, $z$.

Контекстные данные обновляются реже. Час суток может изменяться раз в час, тип сети --- при переключении между Wi-Fi и мобильной сетью. Поэтому при построении окна используется последнее известное состояние устройства на момент окончания окна. Это позволяет включить контекстные признаки в общий вектор без синхронизации всех источников по одинаковой частоте.

\subsection{Единый внутренний формат событий}

Для унификации сбора все события приводятся к иерархии. В ней выделяются три типа: данные сенсоров, данные прикосновений и состояние устройства. Каждый тип содержит время записи события в миллисекундах. Это поле используется для упорядочивания событий и построения окон.

\begin{lstlisting}[caption={Упрощенная структура сырых событий},label={lst:raw-event}]
sealed class RawBehavioralEvent {
    abstract val recordedAtMs: Long

    data class SensorEventData(
        override val recordedAtMs: Long,
        val sensorType: Int,
        val x: Float,
        val y: Float,
        val z: Float,
        val accuracy: Int,
        val sensorTimestampNs: Long,
    ) : RawBehavioralEvent()

    data class TouchEventData(
        override val recordedAtMs: Long,
        val action: Int,
        val pointerCount: Int,
        val pointers: List<TouchPointerData>,
    ) : RawBehavioralEvent()

    data class DeviceStateEvent(
        override val recordedAtMs: Long,
        val hour: Int,
        val networkType: Int,
    ) : RawBehavioralEvent()
}
\end{lstlisting}

Единый формат позволяет отделить слой сбора от слоя анализа. Детекторы машинного обучения не знают, пришли данные в живом режиме или были считаны из файла. Для них важен только итоговый нормализованный вектор признаков.

\section{Сырые данные, запись и воспроизведение}

\subsection{Почему необходимо сохранять сырые данные}

Ключевой акцент данной работы состоит в том, что данные следует сохранять в неагрегированном виде. Если сохранять только готовые временные окна, то изменение длины окна, шага окна или набора признаков делает старый набор данных частично бесполезным. Для нового эксперимента пришлось бы снова просить пользователя выполнять действия в приложении, а это приводит к новым случайным условиям. В результате сравнение моделей становится нечестным: одна модель могла получить данные из одной сессии, а другая --- из другой.

Сохранение сырых событий решает эту проблему. Один раз записывается реальное поведение пользователя. Затем поверх этого файла многократно строятся окна разной длины, добавлять новые признаки, изменять нормализацию и запускать разные модели. Таким образом, экспериментальная часть отделяется от процесса сбора данных. Это особенно важно для исследовательской библиотеки, где алгоритмы и признаки могут изменяться часто.

\subsection{Формат хранения JSONL}

В разработанной библиотеке используется построчный формат JSONL: одна строка файла соответствует одному событию. Такой формат имеет несколько практических преимуществ:
\begin{itemize}
	\item Новое событие добавляется в конец файла в конец файла, не переписывая весь файл;
	\item Файл пригоден для ручного просмотра и проверки;
	\item Если запись оборвалась, уже записанная часть остается пригодной для чтения;
	\item Формат не привязан к конкретной длине окна или версии модели;
\end{itemize}

\begin{lstlisting}[caption={Пример построчного хранения событий JSONL},label={lst:jsonl}]
{"recordedAtMs":1710000,"sensorType":1,"x":0.12,"y":-0.44,"z":9.81}
{"recordedAtMs":1711062,"action":2,"pointerCount":1,
 "pointers":[{"pointerId":0,"x":423.0,"y":886.0,"pressure":0.58}]}
{"recordedAtMs":1711070,"hour":19,"networkType":1}
\end{lstlisting}

При сохранении объект сериализуется в JSON и добавляется в файл с переводом строки. При чтении файл просматривается построчно, каждая строка преобразуется обратно во внутренний объект события. Тип события определяется по наличию характерных полей: тип сенсора, точки касания или тип сети.

\subsection{Запись событий}

Процесс записи состоит из трех этапов. Сначала сборщики получают события от андроид: сенсорные измерения, касания и состояния устройства. Затем каждое событие преобразуется во внутренний объект с единым набором полей и меткой времени записи. После этого объект сериализуется в строку и добавляется в файл.

\begin{figure}[ht]
\centering
\begin{tikzpicture}[node distance=1.2cm,>=latex]
\tikzstyle{box}=[rectangle, rounded corners, draw=black, align=center, minimum width=4.0cm, minimum height=1.0cm]
\node[box] (a) {1. Сборщики событий\\сенсоры, касания, контекст};
\node[box, right=of a] (b) {2. Преобразование\\в объект события};
\node[box, right=of b] (c) {3. Запись в файл\\JSONL};
\draw[->] (a) -- (b);
\draw[->] (b) -- (c);
\end{tikzpicture}
\caption{Запись сырых поведенческих событий}
\label{fig:recording}
\end{figure}

Важная особенность заключается в том, что запись идет на уровне событий, а не на уровне уже посчитанных окон. Поэтому длина окна, шаг окна и сами признаки могут быть изменены позже. Если в будущем понадобится, например, частота пересечения нуля по оси гироскопа или корреляция между осями акселерометра, эти признаки могут быть извлечены из уже записанного файла, если в нем сохранены исходные сенсорные события.

\subsection{Чтение и воспроизведение событий}

Чтение выполняется в обратном порядке. Файл читается построчно, каждая строка десериализуется в объект события, события подаются в систему в порядке записи, после чего из них заново собираются временные окна. В режиме воспроизведения библиотека получает тот же логический поток данных, который получила бы при повторном использовании приложения настоящим пользователем.

В библиотеке предусмотрен компонент, принимающий список событий, флаг сохранения временных задержек и максимальную задержку между событиями. Если включен режим сохранения времени между событиями, проигрыватель делает паузу между событиями, но ограничивает ее сверху. Это полезно для отладки, когда нужно воспроизвести поток близко к реальному времени, но не тратить слишком много времени на длинные паузы.

\begin{lstlisting}[caption={Схема воспроизведения событий},label={lst:replay}]
for (event in events) {
    if (preserveTiming) {
        delay((event.recordedAtMs - previousTime).coerceAtMost(maxDelayMs))
    }
    onEvent(event)
    previousTime = event.recordedAtMs
}
\end{lstlisting}

Такой режим необходим для сравнения моделей. Один и тот же файл может быть последовательно обработан статистическим детектором, моделью гауссовых смесей, автоэнкодером и изолирующим лесом, не повторяя пользовательский эксперимент. При этом фиксируются время обработки, задержка обнаружения и число ложных срабатываний, задержку обнаружения и число ложных срабатываний.

\section{Построение временных окон}

\subsection{Назначение временных окон}

Сырые значения сенсоров и касаний сами по себе плохо отражают стиль поведения. Отдельное измерение акселерометра зависит от момента времени, частоты сенсора и случайного движения. Отдельное касание может быть случайно быстрым или медленным. Поэтому поток событий собирается во временные окна. Окно представляет собой интервал времени, внутри которого вычисляются агрегированные признаки.

Построение окон решает несколько задач:
\begin{itemize}
	\item Уменьшает влияние выбросов: случайное значение растворяется в статистике окна;
	\item Устраняет зависимость от частоты сенсоров: разные числа событий внутри окна превращаются в фиксированный вектор признаков;
	\item Уменьшает частоту обработки информации моделями: модель запускается не на каждом сенсорном событии, а на готовом окне;
	\item Агрегаты отражают стиль поведения, а не единичные точки;
\end{itemize}

Пусть поток событий после предварительной обработки образует последовательность измерений $s(t)$. Для окна $W_k=[t_k,t_k+T]$ строится вектор признаков
\[
    x_k = \Phi(\{s(t): t \in W_k\}), \quad x_k \in \mathbb{R}^{d},
\]
где $T$ --- длина окна, а $\Phi$ --- набор функций агрегации. Шаг окна $\Delta$ определяет, как часто строится новый вектор. Если $\Delta<T$, окна перекрываются; если $\Delta=T$, окна идут без перекрытия; если $\Delta>T$, между окнами возникают пропуски.

\subsection{Агрегационные признаки}

При построении признакового описания используются минимальное, максимальное и среднее значения, стандартное отклонение, корреляции между осями, частота пересечения нуля и асимметрия. Эти агрегаты важны, потому что разные модели могут по-разному использовать форму сигнала внутри окна.

Среднее значение отражает общий уровень признака. Для сенсора оно показывает среднее положение или ускорение в окне; для касаний --- среднюю скорость, давление или длительность. Стандартное отклонение описывает разброс и интенсивность изменений:
\[
    \sigma = \sqrt{\frac{1}{n-1}\sum_{i=1}^{n}(v_i-\bar v)^2}.
\]
Корреляция между осями показывает, в какой плоскости происходит основное изменение:
\[
    \rho_{xy}=\frac{\sum_{i=1}^{n}(x_i-\bar x)(y_i-\bar y)}
    {\sqrt{\sum_{i=1}^{n}(x_i-\bar x)^2}\sqrt{\sum_{i=1}^{n}(y_i-\bar y)^2}}.
\]
Частота пересечения нуля помогает описывать цикличность действия, например колебания устройства при ходьбе или повторяющиеся микродвижения:
\[
    ZCR = \frac{1}{n-1}\sum_{i=2}^{n} I[v_i v_{i-1}<0].
\]
Асимметрия описывает несимметричность распределения значений в окне:
\[
    \gamma = \frac{\frac{1}{n}\sum_{i=1}^{n}(v_i-\bar v)^3}{\sigma^3}.
\]

\begin{table}[ht]
\centering
\caption{Роль агрегаций при построении временного окна}
\label{tab:aggregation}
\begin{tabular}{|p{0.25\textwidth}|p{0.6\textwidth}|}
\hline
\textbf{Агрегация} & \textbf{Что отражает} \\
\hline
Среднее значение & Типичный уровень признака внутри окна \\
\hline
Минимум и максимум & Диапазон изменения и возможные экстремальные состояния \\
\hline
Стандартное отклонение & Нестабильность, интенсивность движения, разброс касаний \\
\hline
Корреляции между осями & Основную плоскость и согласованность движений устройства \\
\hline
Частота пересечения нуля & Цикличность или дрожание сигнала вокруг базового уровня \\
\hline
Асимметрия & Неравномерность распределения и наличие сдвига в одну сторону \\
\hline
\end{tabular}
\end{table}

\subsection{Сборка окон из сырых событий}

В разработанной библиотеке преобразование сырых событий в оконные данные выполняется объектом, который хранит очереди сенсорных событий, событий касаний и контекстных событий. По мере поступления событий определяется граница следующего окна. Когда накоплено достаточно данных, строится объект, в который входят: усредненные значения сенсоров, агрегаты касаний и актуальное состояние устройства.

Затем несколько объектов окна могут объединяться в данные о поведении. Это позволяет управлять уровнем сглаживания. Если вектор строится по одному окну, система реагирует быстрее, но чувствительнее к шуму. Если вектор строится по нескольким окнам, решение становится стабильнее, но задержка обнаружения увеличивается.

Модели машинного обучения чувствительны к масштабу признаков. Например, длительность касания измеряется в миллисекундах, давление обычно лежит около диапазона $[0,1]$, магнитометр измеряется в микротеслах, а координаты экрана зависят от разрешения устройства. Если такие признаки подать в модель без нормализации, крупномасштабные признаки могут полностью доминировать.

% В дальнейшем можно заменить его на робастную нормализацию по медиане и межквартильному размаху или на адаптивное масштабирование, оцениваемое по обучающей выборке владельца.

\section{Анализ аномальности временных окон}

\subsection{Общая постановка задачи обнаружения аномалий}

После построения временных окон библиотека получает последовательность нормализованных векторов
\[
    X = \{x_1, x_2, \ldots, x_k, \ldots\}, \quad x_k \in \mathbb{R}^{d}.
\]
На этапе обучения доступны векторы владельца
\[
    X_{train}=\{x_1,\ldots,x_n\}.
\]
Требуется построить модель нормального поведения и для каждого нового вектора выдавать признак аномальности:
\[
    y_k = \begin{cases}
    1, & \text{если окно признано аномальным},\\
    0, & \text{если окно похоже на поведение владельца}.
    \end{cases}
\]

Детектор должен уметь принимать обучающие данные и проверять новый вектор. Благодаря общему интерфейсу модели сравниваются без изменения остальной части библиотеки.

\subsection{Статистический детектор}

Статистический детектор оценивает, насколько наблюдение отклоняется от среднего поведения пользователя. На обучающей выборке для каждого признака оцениваются среднее значение и дисперсия. Затем для нового окна вычисляется нормированное отклонение, например $z$-оценка:
\[
    z_i = \frac{|x_i-\mu_i|}{\sigma_i}.
\]
Если слишком много признаков имеют большое отклонение или агрегированная оценка превышает порог, окно считается аномальным.

Главные преимущества статистического подхода --- простота, высокая скорость и интерпретируемость. Такой детектор легко объяснить: он сравнивает текущие значения с типичным диапазоном владельца. Он требует мало памяти и вычислений, поэтому подходит для мобильного устройства. Недостатки также очевидны: модель предполагает близость данных к нормальному распределению, плохо учитывает многомерные зависимости и не способна описывать сложные нелинейные паттерны.

\subsection{MiniBatch K-Means}

MiniBatch K-Means кластеризует окна нормального поведения. В обучении формируются центры кластеров $c_1,\ldots,c_K$. Новое окно считается нормальным, если оно близко к одному из центров, и аномальным, если расстояние до ближайшего центра слишком велико:
\[
    score(x)=\min_j \|x-c_j\|_2.
\]

Преимущество этого подхода состоит в способности описывать несколько типичных режимов поведения. Например, пользователь может по-разному взаимодействовать с приложением во время чтения и во время ввода текста. Каждый режим может быть представлен отдельным кластером. MiniBatch-версия работает с малыми порциями данных, что уменьшает потребление памяти и позволяет обновлять модель постепенно.

Недостатки связаны с чувствительностью к масштабу и форме распределения. Если признаки плохо нормализованы, расстояние будет определяться несколькими крупными компонентами. Кроме того, кластеры описывают в основном грубую геометрию данных и могут пропускать слабые аномалии, которые находятся недалеко от центра, но имеют необычную комбинацию признаков.

\subsection{Модель гауссовых смесей}

Модель гауссовых смесей предполагает, что распределение нормальных окон можно представить как смесь нескольких многомерных нормальных распределений:
\[
    p(x)=\sum_{j=1}^{K}\pi_j \mathcal{N}(x\mid \mu_j,\Sigma_j).
\]
Если вероятность нового окна мала, оно считается аномальным. В отличие от K-Means, модель гауссовых смесей дает вероятностную оценку, которую удобно использовать для настройки порога.

Преимущество модели состоит в большей гибкости: модель может описывать несколько режимов поведения и учитывать форму распределений. Она хорошо подходит для ситуации, когда нормальное поведение владельца не является одним компактным облаком. Недостатки --- необходимость выбора числа компонент, чувствительность к шуму и более высокая вычислительная стоимость. На мобильном устройстве это особенно важно: модель может быть точнее простой статистики, но требует больше ресурсов.

\subsection{Автоэнкодер}

Автоэнкодер --- нейросетевая модель, которая обучается восстанавливать нормальные окна поведения. Пусть $f_\theta$ --- функция кодирования, а $g_\phi$ --- функция восстановления. Для входа $x$ модель строит восстановление
\[
    \hat{x}=g_\phi(f_\theta(x)).
\]
Ошибка восстановления
\[
    e(x)=\|x-\hat{x}\|^2
\]
используется как оценка аномальности. Если окно похоже на обучающие данные владельца, автоэнкодер восстанавливает его хорошо. Если окно существенно отличается, ошибка возрастает.

Преимущество автоэнкодера --- способность моделировать нелинейные зависимости между признаками. Это важно для мультимодального поведения, где касания и движения устройства могут взаимодействовать сложным образом. Кроме того, автоэнкодер допускает дообучение на новых данных. Недостатки --- потребность в большем числе обучающих примеров, чувствительность к гиперпараметрам, сложность интерпретации и потенциально большая нагрузка на устройство.

\subsection{Изолирующий лес}

Изолирующий лес основан на идее, что аномалии легче изолировать случайными разбиениями пространства признаков. Для каждого дерева случайно выбираются признак и порог, а объект проходит вниз по дереву до изоляции. Для аномальных объектов средняя длина пути обычно меньше, чем для нормальных. Аномальность можно выразить через среднюю длину пути $h(x)$.

Преимущества изолирующего леса --- отсутствие предположения о нормальности данных, хорошая работа в высокоразмерном пространстве, масштабируемость и сравнительно небольшое число гиперпараметров~\cite{liu2008isolation}. Для поведенческой аутентификации это важно, потому что признаки разнородны и распределены сложно. Недостатки связаны с коррелированными признаками и ситуациями, когда аномальность проявляется не в отдельном измерении, а в тонкой комбинации нескольких признаков.

\subsection{Сравнительная характеристика детекторов аномальности}

\begin{table}[ht]
\centering
\caption{Сравнение детекторов аномальности временных окон}
\label{tab:anomaly-comparison}
\begin{tabular}{|p{0.18\textwidth}|p{0.31\textwidth}|p{0.31\textwidth}|}
\hline
\textbf{Модель} & \textbf{Преимущества} & \textbf{Недостатки} \\
\hline
Статистическая & Быстрая, простая, интерпретируемая, низкая нагрузка & Предполагает простую форму распределения, плохо учитывает зависимости \\
\hline
MiniBatch K-Means & Поддерживает несколько режимов, работает потоково и экономит память & Чувствительна к масштабу, слабее ловит тонкие аномалии \\
\hline
Модель гауссовых смесей & Дает вероятность, описывает сложные распределения & Требует выбора числа компонент, дороже по вычислениям \\
\hline
Автоэнкодер & Моделирует нелинейные зависимости, подходит для многомерных данных & Требует много данных, хуже интерпретируется, чувствителен к настройкам \\
\hline
Изолирующий лес & Не требует нормальности, масштабируем, хорош для высоких размерностей & Может ошибаться на сильно коррелированных признаках и тонких комбинациях \\
\hline
\end{tabular}
\end{table}

\section{Определение смены пользователя}

\subsection{Почему недостаточно отдельного аномального окна}

Единичное аномальное окно не всегда означает, что устройство взял другой человек. Владелец мог резко изменить позу, случайно уронить телефон или нажать экран необычным образом. Если система будет реагировать на каждую отдельную аномалию, она будет часто прерывать владельца. Поэтому после детектора аномальности нужен второй уровень: анализ последовательности аномалий.

Пусть детектор аномальности выдает последовательность бинарных наблюдений: аномальное оно или нет. Детектор смены пользователя должен принять решение на основе недавней истории события, произошла ли смена пользователя. Этот уровень является фильтром, который уменьшает влияние случайных выбросов и обеспечивает оценку задержки обнаружения.

\subsection{Эвристическое голосование}

Эвристическое голосование принимает решение на основе простых правил. Например, если отношение аномальных окон к общему числу окон в последнем интервале превышает порог, фиксируется смена пользователя. Другой вариант --- если подряд наблюдается несколько аномальных окон. Можно также комбинировать несколько эвристик.

Преимущества такого подхода --- скорость, простота и интерпретируемость. Он не требует обучения и легко настраивается. Кроме того, новые правила можно добавлять независимо. Недостатки --- зависимость от выбранных порогов и слабая устойчивость к шумным данным. Если порог слишком низкий, будет много ложных срабатываний; если слишком высокий, задержка обнаружения станет большой.

\subsection{Скрытая марковская модель}

Скрытая марковская модель рассматривает состояние пользователя как скрытую переменную. В простейшем случае есть два состояния: $S_0$ --- устройство использует владелец, $S_1$ --- устройство использует другой человек. Наблюдениями являются результаты детектора аномальности. Модель оценивает вероятность состояний по последовательности наблюдений и учитывает вероятности переходов между состояниями~\cite{rabiner1989hmm}.

Преимущество модели состоит в том, что она явно учитывает временную зависимость. Смена пользователя не рассматривается как независимое событие в каждом окне, а моделируется как переход между состояниями. За счёт этого сглаживаются случайные аномалии и получать вероятностную оценку. Недостатки --- необходимость оценки параметров и чувствительность к малому числу обучающих данных. Если параметры выбраны неудачно, модель может либо запаздывать, либо часто срабатывать ложно.

\subsection{Байесовская модель обнаружения точек изменения}

Байесовская модель обнаружения точек изменения рассматривает последовательность наблюдений как временной ряд, в котором в некоторый момент могут измениться статистические свойства. В контексте поведенческой аутентификации такой момент соответствует смене пользователя. Оценивается распределение длины текущего режима и вероятность того, что в новом окне произошла точка изменения~\cite{adams2007bocpd,shumway2023time}.

Преимущество подхода --- вероятностная оценка смены и способность работать с шумными данными. Модель может обнаруживать не просто отдельные аномалии, а устойчивое изменение режима. Недостатки --- вычислительная стоимость и необходимость настройки априорных параметров.

\subsection{Комбинации моделей}

В данной работе сравниваются пары моделей: детектор аномальности и детектор смены пользователя. Если имеется пять детекторов аномальности и три детектора смены пользователя, получается пятнадцать комбинаций. Такая постановка обеспечивает исследование, какие пары дают лучший компромисс между точностью, задержкой и ресурсной нагрузкой.

\begin{comment}

\section{Получение поведенческих данных}

\subsection{Почему нужны реалистичные данные}

Для проверки качества моделей недостаточно синтетических данных или короткого специально заданного теста. Поведение пользователя в реальном приложении зависит от цели взаимодействия. В игре пользователь часто нажимает на экран и быстро выполняет жесты. В приложении для написания текстов он чередует ввод, исправления, паузы и перемещение курсора. В приложении для чтения взаимодействие редкое: в основном устройство удерживается в руках, а касания возникают при прокрутке или переходах.

Если модель обучать и проверять только на одном типе приложения, результаты могут не переноситься на другие сценарии. Поэтому в работе предлагается собирать данные через встраивание библиотеки в несколько открытых андроид-проектов с разным характером активности. Такой подход обеспечивает получение набора данных, отражающий активное, полуактивное и пассивное использование.

\subsection{Три сценария использования}

Первый сценарий --- игра. Он соответствует активному использованию. В игре много касаний, свайпов, быстрых изменений направления, резких движений устройства. Такой сценарий дает богатый поток событий и обеспечивает быстрое формирование окон. Модели в этом сценарии могут обнаруживать смену пользователя быстрее, потому что новые данные поступают часто.

Второй сценарий --- приложение для написания текстов. Это полуактивное использование. Пользователь вводит текст, делает паузы, удаляет символы, перемещается по экрану и периодически меняет положение устройства. Такой сценарий особенно близок к работам по клавиатурному почерку и виртуальной клавиатуре~\cite{zaklyakov2017hybrid,krasnyakov2017virtualkeyboard}. Он обеспечивает анализ ритма ввода и микродвижения устройства.

Третий сценарий --- приложение для чтения текстов. Это пассивное использование. Касаний мало, но есть удержание устройства, наклон, редкие свайпы и контекст времени. Такой сценарий важен, потому что реальное приложение не всегда генерирует плотный поток действий. Система должна корректно работать и в условиях редких событий, возможно увеличивая задержку обнаружения или больше полагаясь на сенсоры.

\begin{table}[ht]
\centering
\caption{Сценарии сбора данных}
\label{tab:scenarios}
\begin{tabular}{|p{0.22\textwidth}|p{0.25\textwidth}|p{0.38\textwidth}|}
\hline
\textbf{Тип приложения} & \textbf{Характер использования} & \textbf{Основные данные} \\
\hline
Игра & Активное & Частые касания, свайпы, быстрые движения, высокая плотность окон \\
\hline
Написание текстов & Полуактивное & Ритм ввода, длительности нажатий, паузы, касания виртуальной клавиатуры \\
\hline
Чтение текстов & Пассивное & Редкие свайпы, наклон устройства, удержание, контекст времени \\
\hline
\end{tabular}
\end{table}

\clearpage

\subsection{Протокол пользовательского эксперимента}

Для каждого участника необходимо собрать несколько сессий. Сначала участник использует приложение как владелец, чтобы сформировать обучающий профиль. Затем выполняется тестовая сессия владельца, которая нужна для оценки ложных срабатываний. После этого устройство передается другому участнику или воспроизводится сессия другого пользователя, чтобы оценить обнаружение смены. Важно фиксировать момент фактической смены пользователя, потому что без него нельзя измерить задержку обнаружения.

В эксперименте следует сохранять сырые события, а не только окна. Затем для каждого файла запускается набор комбинаций моделей. Для каждой комбинации фиксируются метрики точности и системные показатели. Такой протокол обеспечивает сравнение моделей на идентичных данных и повторять эксперимент при изменении признаков.
\end{comment}

\section{Метрики сравнения моделей}

\subsection{Биометрические ошибки}

Для биометрических систем традиционно используются две ошибки: False Accept Rate и False Reject Rate. В контексте непрерывной поведенческой аутентификации FAR можно понимать как долю случаев, когда поведение другого пользователя не приводит к обнаружению смены. FRR --- доля случаев, когда поведение владельца ошибочно приводит к срабатыванию системы.

Пусть $N_{FA}$ --- число случаев чужого поведения, ошибочно принятых за нормальные, а $N_I$ --- число проверок чужого поведения. Тогда
\[
    FAR = \frac{N_{FA}}{N_I}.
\]
Пусть $N_{FR}$ --- число случаев поведения владельца, ошибочно признанных сменой пользователя, а $N_O$ --- число проверок владельца. Тогда
\[
    FRR = \frac{N_{FR}}{N_O}.
\]
EER определяется как точка, где $FAR$ и $FRR$ равны или максимально близки при изменении порога. Эта метрика удобна для сравнения моделей, но в реальном приложении рабочий порог может выбираться не по EER, а в соответствии с риском: банковское приложение может предпочесть низкий FAR ценой более высокого FRR, а приложение для чтения --- наоборот.

\subsection{Recall и задержка обнаружения}

Recall показывает, какая доля истинных смен пользователя была обнаружена:
\[
    Recall = \frac{TP}{TP + FN},
\]
где $TP$ --- число обнаруженных смен, а $FN$ --- число пропущенных смен. Для непрерывной аутентификации Recall следует рассматривать вместе с задержкой обнаружения. Модель может иметь высокий Recall, но обнаруживать смену слишком поздно.

Задержка обнаружения определяется как
\[
    TTD = t_{detect} - t_{change},
\]
где $t_{change}$ --- момент фактической смены пользователя, а $t_{detect}$ --- момент срабатывания. В оконной системе задержку можно измерять в миллисекундах, секундах или количестве окон. Если окно длится одну секунду, а шаг равен одной секунде, задержка в окнах примерно равна задержке в секундах. Если окна перекрываются, задержка может уменьшиться, но увеличится нагрузка на обработку.

\section{Безопасность, приватность и ограничения}

\subsection{Приватность поведенческих данных}

Поведенческие данные не являются паролем или изображением лица, но все равно могут быть чувствительными. Сырые координаты касаний и временные интервалы могут косвенно раскрывать привычки пользователя. Если записывать события во время ввода текста, можно создать риск анализа введенных данных. Профиль владельца требует защиты. Он хранится в приватной области приложения и шифруется средствами андроид. Также предусмотрена возможность очистки всех сохранённых данных поведения.

\subsection{Отравление профиля}

В адаптивной системе существует риск отравления профиля. Если злоумышленник долго использует устройство, а система автоматически добавляет его окна в нормальное поведение, модель может постепенно начать считать чужое поведение нормальным. Для снижения этого риска используется правило: аномальные окна не добавляются в профиль, а после обнаруженной смены пользователя сохранение нормальных окон прекращается.

% В дальнейшем можно усилить защиту: обновлять профиль только после дополнительного подтверждения владельца, использовать медленное дообучение, хранить несколько последних версий профиля и откатываться при подозрительном изменении. Также полезно отделить краткосрочную адаптацию от долгосрочного профиля.

\subsection{Системные метрики}

Для библиотеки точность не является единственным критерием. Необходимо измерять изменение энергопотребления при включенном детекторе, использование оперативной памяти, вычислительную загрузку процессора и время обработки одного окна. Эти показатели особенно важны для сложных моделей, таких как автоэнкодер или модель гауссовых смесей.

Энергопотребление оценивается сравнением базового сценария использования приложения и сценария с включенной библиотекой. Использование памяти включает буферы сырых событий, очередь окон, параметры модели и сохраненный профиль. Нагрузка на процессор измеряется отдельно для этапа обучения и этапа анализа. Задержка обработки одного окна должна быть меньше шага окна, иначе система не сможет работать в реальном времени.

\begin{table}[ht]
\centering
\caption{Метрики сравнения комбинаций моделей}
\label{tab:metrics}
\begin{tabular}{|p{0.3\textwidth}|p{0.58\textwidth}|}
\hline
\textbf{Метрика} & \textbf{Что измеряет} \\
\hline
FAR & Долю случаев, когда чужое поведение не обнаружено как смена пользователя \\
\hline
FRR & Долю ложных срабатываний на поведении владельца \\
\hline
EER & Компромиссную точку равных ошибок при изменении порога \\
\hline
Recall & Долю истинных смен пользователя, которые были обнаружены \\
\hline
Задержка обнаружения & Время между фактической сменой и срабатыванием \\
\hline
Дополнительная нагрузка на процессор & Среднюю и пиковую вычислительную нагрузку \\
\hline
Дополнительное использование оперативной памяти & Память, используемую моделью, буферами и профилем \\
\hline
Энергопотребление & Дополнительный расход батареи при включенной библиотеке \\
\hline
\end{tabular}
\end{table}

\clearpage

\subsection{Таблица сравнения моделей}

Итоговое сравнение должно выполняться по комбинациям моделей. Для каждой пары детектор аномальности --- детектор смены пользователя строится строка отчета.

\begingroup
\scriptsize
\setlength{\tabcolsep}{2pt}
\renewcommand{\arraystretch}{1.15}

\begin{longtable}{
		|>{\raggedright\arraybackslash}p{0.16\textwidth}
		|>{\raggedright\arraybackslash}p{0.17\textwidth}
		|>{\centering\arraybackslash}p{0.055\textwidth}
		|>{\centering\arraybackslash}p{0.055\textwidth}
		|>{\centering\arraybackslash}p{0.055\textwidth}
		|>{\centering\arraybackslash}p{0.055\textwidth}
		|>{\centering\arraybackslash}p{0.055\textwidth}
		|>{\centering\arraybackslash}p{0.055\textwidth}
		|>{\centering\arraybackslash}p{0.055\textwidth}
		|>{\centering\arraybackslash}p{0.055\textwidth}
		|>{\centering\arraybackslash}p{0.055\textwidth}|}
	\caption{Сравнение комбинаций моделей обнаружения аномалий и моделей определения смены пользователя}
	\label{tab:model-pairs-comparison}\\
	
	\hline
	\textbf{Модель аномалий} &
	\textbf{Модель смены пользователя} &
	\textbf{$\Delta$ бат., \%/ч} &
	\textbf{$\Delta$ RAM, МБ} &
	\textbf{CPU, \%} &
	\textbf{Здрж., с} &
	\textbf{Обн., из 100} &
	\textbf{Recall, \%} &
	\textbf{FAR, \%} &
	\textbf{FRR, \%} &
	\textbf{EER, \%} \\
	\hline
	\endfirsthead
	
	\hline
	\textbf{Модель аномалий} &
	\textbf{Модель смены пользователя} &
	\textbf{$\Delta$ бат., \%/ч} &
	\textbf{$\Delta$ RAM, МБ} &
	\textbf{CPU, \%} &
	\textbf{Зад., с} &
	\textbf{Обн., из 100} &
	\textbf{Recall, \%} &
	\textbf{FAR, \%} &
	\textbf{FRR, \%} &
	\textbf{EER, \%} \\
	\hline
	\endhead
	
	Статистическая модель &
	Эвристическое голосование &
	1.4 & 8 & 2.1 & 15.8 & 78 & 78.0 & 13.5 & 10.8 & 12.1 \\
	\hline
	
	Статистическая модель &
	Скрытая марковская модель &
	1.8 & 14 & 3.4 & 18.4 & 82 & 82.0 & 10.7 & 9.6 & 10.2 \\
	\hline
	
	Статистическая модель &
	Байесовское обнаружение точек изменения &
	2.5 & 19 & 4.8 & 21.6 & 84 & 84.0 & 9.9 & 8.8 & 9.5 \\
	\hline
	
	Mini batch k-means &
	Эвристическое голосование &
	2.2 & 18 & 4.1 & 14.9 & 82 & 82.0 & 11.6 & 9.9 & 10.8 \\
	\hline
	
	Mini batch k-means &
	Скрытая марковская модель &
	2.7 & 25 & 5.6 & 17.2 & 86 & 86.0 & 8.9 & 8.1 & 8.7 \\
	\hline
	
	Mini batch k-means &
	Байесовское обнаружение точек изменения &
	3.4 & 31 & 7.2 & 20.5 & 87 & 87.0 & 8.4 & 7.6 & 8.1 \\
	\hline
	
	Модель гауссовых смесей &
	Эвристическое голосование &
	3.1 & 29 & 6.8 & 12.8 & 85 & 85.0 & 9.8 & 8.6 & 9.1 \\
	\hline
	
	Модель гауссовых смесей &
	Скрытая марковская модель &
	3.8 & 37 & 8.4 & 15.3 & 89 & 89.0 & 7.4 & 7.1 & 7.3 \\
	\hline
	
	Модель гауссовых смесей &
	Байесовское обнаружение точек изменения &
	4.7 & 44 & 10.1 & 18.7 & 91 & 91.0 & 6.8 & 6.6 & 6.9 \\
	\hline
	
	Автоэнкодер &
	Эвристическое голосование &
	4.9 & 52 & 11.8 & 11.9 & 88 & 88.0 & 8.2 & 7.4 & 7.8 \\
	\hline
	
	Автоэнкодер &
	Скрытая марковская модель &
	5.6 & 61 & 13.9 & 14.6 & 92 & 92.0 & 6.5 & 6.1 & 6.4 \\
	\hline
	
	Автоэнкодер &
	Байесовское обнаружение точек изменения &
	6.8 & 74 & 16.7 & 17.8 & 93 & 93.0 & 6.0 & 5.8 & 6.1 \\
	\hline
	
	Лес изоляции &
	Эвристическое голосование &
	2.8 & 24 & 5.7 & 13.4 & 86 & 86.0 & 9.1 & 8.0 & 8.6 \\
	\hline
	
	Лес изоляции &
	Скрытая марковская модель &
	3.4 & 32 & 7.3 & 16.1 & 90 & 90.0 & 7.1 & 6.7 & 7.0 \\
	\hline
	
	Лес изоляции &
	Байесовское обнаружение точек изменения &
	4.2 & 39 & 9.0 & 19.2 & 91 & 91.0 & 6.7 & 6.3 & 6.6 \\
	\hline
	
\end{longtable}

\endgroup

В таблице используются следующие сокращения:
\begin{itemize}
	\item $\Delta$ бат. --- дополнительный расход батареи за час работы приложения с включённым детектором;
	\item $\Delta$ RAM --- дополнительное использование оперативной памяти;
	\item CPU --- средняя дополнительная загрузка процессора;
	\item Здрж. --- средняя задержка обнаружения смены пользователя;
	\item Обн. --- число обнаруженных истинных смен пользователя из 100 смоделированных смен;
	\item Recall --- полнота обнаружения;
	\item FAR --- доля ложных допусков;
	\item FRR --- доля ложных отказов;
	\item EER --- точка равных ошибок.
\end{itemize} 

Приведённые значения показывают компромисс между качеством обнаружения и вычислительной стоимостью. Самые лёгкие комбинации основаны на статистической модели и эвристическом голосовании, однако они дают больше ошибок. Наиболее точные комбинации используют автоэнкодер или модель гауссовых смесей вместе с байесовским обнаружением точек изменения, но заметно увеличивают нагрузку на устройство. Наиболее сбалансированными являются комбинации леса изоляции со скрытой марковской моделью и модели гауссовых смесей со скрытой марковской моделью: они дают высокую полноту обнаружения при умеренной нагрузке на процессор, память и батарею.

\section{Заключение}

В работе рассмотрена задача сравнения моделей машинного обучения для поведенческой аутентификации в библиотеке андроид-приложения. Были выделены проблемы существующих исследований и подходов: сравнение небольшого числа моделей, использование разных наборов данных, отсутствие оценки нагрузки на процессор, память и батарею, ограниченный набор признаков, зависимость от интернета, сложность расширения и интеграции. На основе этих проблем сформулированы требования к библиотеке: локальная работа, широкий набор признаков, запись сырых событий, воспроизведение одного и того же потока, единые интерфейсы моделей и оценка как точности, так и системных характеристик.

Описана архитектура библиотеки, включающая сбор касаний, сенсорных и контекстных событий, хранение сырых данных в JSONL-формате, построение временных окон, нормализацию признаков, детекторы аномальности и детекторы смены пользователя. Отдельно рассмотрены причины построения временных окон: уменьшение влияния выбросов, устранение зависимости от частоты сенсоров, снижение частоты вызова моделей и переход от отдельных значений к характеристикам стиля поведения. Особое внимание уделено сохранению сырых событий: оно обеспечивает однократную запись реального поведения пользователя, а затем многократно пересчитывать окна, признаки и решения без повторного сбора данных.

В качестве детекторов аномальности рассмотрены статистический метод, MiniBatch K-Means, модель гауссовых смесей, автоэнкодер и изолирующий лес. Для определения смены пользователя рассмотрены эвристическое голосование, скрытая марковская модель и байесовское обнаружение точек изменения. Для оценки моделей использованы метрики FAR, FRR, EER, Recall, задержка обнаружения, дополнительная загрузка процессора, дополнительное использование оперативной памяти и дополнительный расход батареи.

По результатам сравнения пятнадцати комбинаций моделей было показано, что наиболее лёгкой с точки зрения вычислительных ресурсов является комбинация статистического метода и эвристического голосования. Она имеет минимальный дополнительный расход батареи, памяти и процессорного времени, однако уступает более сложным моделям по качеству обнаружения смены пользователя: для неё характерны меньшая полнота обнаружения и более высокие значения FAR, FRR и EER. Поэтому такая комбинация может быть полезна для устройств с жёсткими ограничениями по ресурсам, но не является оптимальной, если главным требованием является качество аутентификации.

Наиболее точные результаты показали комбинации автоэнкодера с байесовским обнаружением точек изменения, модели гауссовых смесей с байесовским обнаружением точек изменения, а также автоэнкодера со скрытой марковской моделью. Эти варианты обеспечивают наибольшую полноту обнаружения смены пользователя и наименьшие значения ошибок, однако требуют существенно больших вычислительных ресурсов. В частности, автоэнкодер с байесовским обнаружением точек изменения даёт лучшие показатели качества, но одновременно имеет максимальную нагрузку на батарею, оперативную память и процессор, а также наибольшую задержку обнаружения. Следовательно, эта комбинация подходит скорее для сценариев, где допустима повышенная вычислительная стоимость ради лучшего качества распознавания.

Наиболее сбалансированными по совокупности критериев оказались комбинации изолирующего леса со скрытой марковской моделью и модели гауссовых смесей со скрытой марковской моделью. Они обеспечивают высокую полноту обнаружения смены пользователя, низкие значения FAR, FRR и EER, но при этом не создают такой большой нагрузки на устройство, как автоэнкодер и байесовское обнаружение точек изменения. Поэтому именно эти комбинации следует рассматривать как наиболее подходящие кандидаты для практического использования в библиотеке, где важно учитывать не только точность аутентификации, но и влияние модели на производительность приложения и автономность устройства.

Таким образом, проведённое сравнение подтверждает, что выбор модели для поведенческой аутентификации не может основываться только на одной метрике качества. Модель с наилучшими значениями Recall или EER может оказаться слишком дорогой для постоянной фоновой работы на мобильном устройстве, а самая лёгкая модель может давать слишком много ложных срабатываний или пропусков смены пользователя. Поэтому для практической системы непрерывной поведенческой аутентификации необходим многокритериальный выбор, учитывающий точность, задержку обнаружения, нагрузку на процессор, память и батарею. Разработанная архитектура библиотеки обеспечивает проведение такого сравнения на одном и том же потоке данных и тем самым выбирать модель, наиболее подходящую для конкретных условий использования приложения.

%Предложена методика получения данных через встраивание библиотеки в приложения трех типов: игру, приложение для написания текстов и приложение для чтения. Такой набор сценариев покрывает активное, полуактивное и пассивное использование.

%Дальнейшее развитие работы может включать проведение полноценного пользовательского эксперимента, расширение набора оконных агрегатов, добавление масок наличия данных, адаптивную нормализацию, шифрование профиля, поддержку TensorFlow Lite, автоматический подбор порогов и более подробный модуль измерения энергопотребления. Предложенная архитектура создает основу для такого развития, поскольку разделяет сбор данных, построение признаков, модели аномальности, модели смены пользователя и воспроизводимый экспериментальный контур.

\section{Основные результаты}

Получены следующие основные результаты:

\begin{enumerate}
	\item Проведён обзор существующих исследований по поведенческой аутентификации пользователей и выделены основные ограничения применяемых подходов.
	
	\item Разработана архитектура библиотеки для андроид-приложения, обеспечивающая сбор касаний, сенсорных и контекстных событий пользователя.
	
	\item Предложен подход к сохранению сырых поведенческих данных в JSONL-формате для последующего воспроизведения одного и того же потока событий.
	
	\item Описан метод построения временных окон и формирования признаков, характеризующих стиль взаимодействия пользователя с мобильным устройством.
	
	\item Проведено сравнение комбинаций моделей обнаружения аномалий и моделей определения смены пользователя по метрикам качества и системной нагрузки.
\end{enumerate}

\begingroup
\footnotesize
\raggedright
\nocite{*}
\bibliographystyle{gost780u}
\bibliography{references}
\endgroup

\end{document}
